\subsection{Intrinsic native model ideal functions from genuine ideal-sheaf data}
\label{subsec:native-model-ideals}

Mathlib's actual $\operatorname{IdealSheafData}$ specifies an ideal on
\emph{every} genuine affine open and compatibility with genuine basic-open
localization. It is the existing scheme ideal-sheaf data object, not an
assumption that the desired function exists. Mathlib has not yet packaged
these data as an actual ideal subsheaf; no additional such packaging or
analytic coherent-sheaf gluing is asserted here. The algebraic stalk
construction works on every scheme. Its model-function application uses
an actual proper structure morphism over an arbitrary prescribed seminorm
bounded on all coefficients. No Noetherian or global finite-generation
hypothesis is imposed. Finite-generator formulas retain the explicit
local generating-family hypothesis.

\begin{definition}[Intrinsic ideal of genuine section germs]
\label{scheme:ideal-sheaf-data-stalk}
\lean{AlgebraicGeometry.Scheme.IdealSheafData.stalkIdeal,
AlgebraicGeometry.Scheme.IdealSheafData.affineGermIdeal_eq_of_le,
AlgebraicGeometry.Scheme.IdealSheafData.affineGermIdeal_eq,
AlgebraicGeometry.Scheme.IdealSheafData.stalkIdeal_eq_affine,
AlgebraicGeometry.Scheme.IdealSheafData.germ_mem_stalkIdeal}
\leanok
Let $I$ be actual ideal-sheaf data on $X$. For $x\in X$ define its stalk
ideal to be the supremum of the genuine germ-map images $I(U)\mathcal O_{X,x}$
over all actual affine opens $U$ containing $x$. For \emph{any} such affine
open, this intrinsic ideal equals that one germ-map image.
\end{definition}
\begin{proof}
For an affine inclusion $W\subset U$, the genuine localization theorem
$I(U)\Gamma(X,W)=I(W)$ and compatibility of actual restriction with germs
show equality of the two germ-map images in the stalk. Given affine
neighborhoods $U,V$ of $x$, the actual affine basis supplies an affine
$W$ containing $x$ inside $U\cap V$. The preceding inclusion equality
therefore makes the images from $U$ and $V$ equal. Every term in the
intrinsic supremum is thus equal to the image from the chosen affine
neighborhood; that image is itself one of the terms, proving both ideal
inclusions. Genuine section germs of $I(U)$ belong to this ideal by
actual ideal-map membership. No stalk ideal equality is assumed.
\end{proof}

\begin{lemma}[Actual stalk ideal arithmetic]
\label{scheme:ideal-sheaf-data-stalk-arithmetic}
\lean{AlgebraicGeometry.Scheme.IdealSheafData.paperStalkIdeal_bot,
AlgebraicGeometry.Scheme.IdealSheafData.paperStalkIdeal_top,
AlgebraicGeometry.Scheme.IdealSheafData.paperStalkIdeal_sup,
AlgebraicGeometry.Scheme.IdealSheafData.paperStalkIdeal_mul,
AlgebraicGeometry.Scheme.IdealSheafData.paperStalkIdeal_pow}
\leanok
\uses{scheme:ideal-sheaf-data-stalk}
The intrinsic stalk ideals preserve zero, unit, ideal sums, products and
all powers, including exponent zero.
\end{lemma}
\begin{proof}
Choose an actual affine neighborhood from the genuine affine basis.
The preceding theorem identifies each intrinsic ideal with its image
from that common neighborhood. Mathlib's actual ideal-sheaf-data arithmetic
is computed on its affine section ideals. The genuine ideal-map laws
preserve their zero, unit, sums, products and powers, proving every
stalk identity. No reducedness or surjectivity of the germ map is used.
\end{proof}

\begin{definition}[Intrinsic actual centered model norm]
\label{scheme:native-model-ideal-function}
\lean{EnvelopeProperty.Berkovich.RelativeSchemeValuedPoint.centerStalkSpectrum,
EnvelopeProperty.Berkovich.RelativeSchemeValuedPoint.isBounded_centerStalkSpectrum,
EnvelopeProperty.Berkovich.RelativeSchemeValuedPoint.centerStalkSpectrum_comap_germ,
EnvelopeProperty.Berkovich.RelativeSchemeValuedPoint.modelIdealNorm,
EnvelopeProperty.Berkovich.RelativeSchemeValuedPoint.modelIdealLog,
EnvelopeProperty.Berkovich.RelativeSchemeValuedPoint.modelIdealNorm_eq_affine,
EnvelopeProperty.Berkovich.RelativeSchemeValuedPoint.modelIdealLog_eq_affine,
EnvelopeProperty.Berkovich.RelativeSchemeValuedPoint.modelIdealNorm_le_one,
EnvelopeProperty.Berkovich.RelativeSchemeValuedPoint.modelIdealLog_le_zero}
\leanok
\uses{scheme:ideal-sheaf-data-stalk,scheme:proper-center-local-bounds,
scheme:relative-section-coefficients,def:general-ideal-norm,lem:general-ideal-pullback}
Let $y$ be a native prescribed-base point of an actual proper scheme,
with true model center $c(y)$. Evaluate the intrinsic stalk ideal of $I$
at the actual bounded center-stalk seminorm to define its norm; its
logarithm defines the model ideal function. For any actual affine center
neighborhood $U$, these equal the ordinary ideal norm and logarithm of
$I(U)$ at the genuine native section-spectrum point. The norm is at most
one and the logarithm at most zero.
\end{definition}
\begin{proof}
The already constructed proper center-stalk valuation gives a genuine
bounded seminorm point of $\mathcal O_{X,c(y)}$. Its pullback along the
actual germ map equals the entire original native section-spectrum
point, because their evaluations agree on every section. The intrinsic
stalk ideal equals the germ-image of $I(U)$ by the first theorem. Ideal
norm pullback at the genuinely bounded target stalk therefore identifies
the intrinsic norm with the actual affine section-ideal norm. Applying
the actual extended logarithm gives the logarithmic identity. Genuine
center-stalk boundedness proves the two bounds. These follow from the
actual proper valuative lift, rather than an assumed coefficient action
or an assumed bound on unrelated section elements.
\end{proof}

\begin{theorem}[Actual affine generators give the same intrinsic function]
\label{scheme:native-model-ideal-generators}
\lean{EnvelopeProperty.Berkovich.RelativeSchemeValuedPoint.modelIdealNorm_eq_generators,
EnvelopeProperty.Berkovich.RelativeSchemeValuedPoint.modelIdeal_generators_eq}
\leanok
\uses{scheme:native-model-ideal-function,lem:general-ideal-norm-generators}
If a finite family $S$ generates the actual ideal $I(U)$ on an affine
center neighborhood, the intrinsic norm equals the maximum of the
original native values of those sections. Any second actual affine center
neighborhood and finite generating family give the same maximum.
The empty generating family computes the genuine zero ideal.
\end{theorem}
\begin{proof}
The previous affine comparison and genuine boundedness of the actual
section-spectrum point allow the proved finite-generator ideal-norm
formula. All section values are original native residue values by the
actual point definition. Both affine maxima equal the one intrinsic
stalk norm, so injectivity of the nonnegative-real inclusion into the
extended nonnegative reals makes them equal. This proves independence
across different affine neighborhoods and their section rings, not only
for two generators of one ideal in one ring. No finite family is assumed
to exist unless explicitly supplied.
\end{proof}

\begin{theorem}[Genuine centered model ideal arithmetic]
\label{scheme:native-model-ideal-arithmetic}
\lean{EnvelopeProperty.Berkovich.RelativeSchemeValuedPoint.modelIdealNorm_bot,
EnvelopeProperty.Berkovich.RelativeSchemeValuedPoint.modelIdealNorm_top,
EnvelopeProperty.Berkovich.RelativeSchemeValuedPoint.modelIdealNorm_sup,
EnvelopeProperty.Berkovich.RelativeSchemeValuedPoint.modelIdealNorm_mul,
EnvelopeProperty.Berkovich.RelativeSchemeValuedPoint.modelIdealNorm_pow,
EnvelopeProperty.Berkovich.RelativeSchemeValuedPoint.modelIdealLog_bot,
EnvelopeProperty.Berkovich.RelativeSchemeValuedPoint.modelIdealLog_top,
EnvelopeProperty.Berkovich.RelativeSchemeValuedPoint.modelIdealLog_sup,
EnvelopeProperty.Berkovich.RelativeSchemeValuedPoint.modelIdealLog_mul,
EnvelopeProperty.Berkovich.RelativeSchemeValuedPoint.modelIdealLog_pow}
\leanok
\uses{scheme:ideal-sheaf-data-stalk-arithmetic,scheme:native-model-ideal-function,
lem:general-ideal-norm-arithmetic,lem:general-ideal-log}
Zero and unit model ideals have norms zero and one and logarithms
$-\infty$ and zero. Sums take maxima, products multiply norms and add
logarithms, and every natural ideal power scales the logarithm by that
exponent. These formulas include zero exponents and negative-infinite
values with the actual extended-real conventions.
\end{theorem}
\begin{proof}
Actual stalk ideal arithmetic identifies the relevant intrinsic ideals.
Apply the established bounded seminorm ideal calculus to the actual
center-stalk point. Unit and zeroth-power identities use its genuine
boundedness. Applying the true extended logarithm gives maxima, products
and powers with their already proved infinity conventions. This constructs
a global intrinsic native model ideal function but does not prove global
continuity, center anti-continuity, analytic sheaf realization or decreasing
function-limit transport. Those need their own explicit producers.
\end{proof}
